\documentclass[10pt,twocolumn]{article}

\usepackage[margin=1in]{geometry}
\usepackage{amsmath}
\usepackage{graphicx}
\usepackage{booktabs}
\usepackage{hyperref}
\usepackage{listings}
\usepackage{xcolor}
\usepackage{caption}
\usepackage{subcaption}
\usepackage{enumitem}

\hypersetup{colorlinks=true, linkcolor=black, citecolor=black, urlcolor=blue}

\lstset{
  basicstyle=\ttfamily\footnotesize,
  breaklines=true,
  frame=single,
  backgroundcolor=\color{gray!10},
}

% -----------------------------------------------------------------------
\title{\textbf{Entropy Source Characterization of the ESP32 Hardware RNG}\\
       \large ECE 580: Hardware Security --- Final Project}
\author{Divija Durga\\Princeton University}
\date{\today}
% -----------------------------------------------------------------------

\begin{document}
\maketitle

% -----------------------------------------------------------------------
\section{Introduction}
% -----------------------------------------------------------------------

Random number generation is a foundational primitive in hardware security.
Cryptographic keys, nonces, and session identifiers derive their security
entirely from the unpredictability of the underlying entropy source.
Weak or biased RNG output has historically led to catastrophic key-recovery
attacks~\cite{debian2008,androidecdsa}.

The ESP32 integrates a hardware true-random-number generator (TRNG) whose
output quality is documented to depend on which peripherals are
active~\cite{esp32trm}.
Specifically, the RF subsystem (Wi-Fi / Bluetooth) and the SAR ADC each
contribute noise to the entropy pool that feeds the hardware RNG register
(\texttt{WDEV\_RND\_REG}).
When neither is active, entropy is derived solely from the internal RC
oscillator, which the manufacturer acknowledges produces lower-quality
randomness.

This work empirically measures the effect of each entropy source
(and their combinations) on RNG output quality using standard statistical
tests and entropy estimation.

% -----------------------------------------------------------------------
\section{Background}
% -----------------------------------------------------------------------

\subsection{ESP32 RNG Architecture}

The ESP32 hardware RNG is a single 32-bit register continuously updated
by hardware noise sources.
The \texttt{esp\_random()} API reads this register directly.
Three classes of entropy sources are documented:

\begin{enumerate}
  \item \textbf{RF subsystem} --- Wi-Fi and Bluetooth transmit/receive
        chains inject broadband thermal noise.
  \item \textbf{SAR ADC} --- Each conversion result contributes shot and
        thermal noise from the analog front-end.
  \item \textbf{Internal RC oscillator} --- Always present; provides
        jitter-based entropy when RF and ADC are inactive.
\end{enumerate}

\subsection{Statistical Tests}

We employ the following metrics and tests.

\paragraph{Shannon entropy.}
\[
  H = -\sum_{i=0}^{255} p_i \log_2 p_i \quad \text{(bits/byte)}
\]
Ideal uniform output achieves $H = 8$.

\paragraph{Min-entropy.}
\[
  H_\infty = -\log_2\!\left(\max_i\, p_i\right)
\]
The security-relevant lower bound: measures the adversary's best
single-guess probability.

\paragraph{NIST Frequency (Monobit) Test~\cite{nist80022}.}
Tests that the proportion of 1-bits is close to $\frac{1}{2}$.
The test statistic is
\[
  s_\text{obs} = \frac{|S_n|}{\sqrt{n}}, \qquad
  p = \mathrm{erfc}\!\left(\frac{s_\text{obs}}{\sqrt{2}}\right).
\]
Pass criterion: $p \geq 0.01$.

\paragraph{NIST Block Frequency Test.}
Divides the bit sequence into $N$ blocks of $M = 128$ bits and tests
that each block has approximately $M/2$ ones.
\[
  \chi^2 = 4M \sum_{i=1}^{N}\!\left(\pi_i - \tfrac{1}{2}\right)^2, \qquad
  p = \Gamma\!\left(\tfrac{N}{2},\, \tfrac{\chi^2}{2}\right).
\]

\paragraph{NIST Runs Test.}
Counts the total number of alternating bit runs and compares to the
expected count for a uniform sequence.

\paragraph{Chi-squared byte uniformity.}
Tests the null hypothesis that all 256 byte values are equally likely
($\chi^2$ with 255 degrees of freedom).

\paragraph{Visual and temporal diagnostics.}
Byte-value histograms and a $256\times256$ bit matrix provide a visual
check for gross non-uniformity or periodic structure.
Consecutive-sample scatter plots ($x_i$ vs.\ $x_{i+1}$) reveal serial
correlation between adjacent outputs.
Byte-level autocorrelation at lags 1--50 tests whether knowledge of
byte $i$ predicts byte $i+k$; for an ideal RNG all coefficients should
lie within the $\pm2/\sqrt{n}$ significance band.

We run all 15 NIST SP 800-22 tests (Table~\ref{tab:nist}).
Tests 7--15 require at least $10^6$ bits; our 4\,MB dataset
($\approx33.5\,\text{Mbit}$) satisfies this requirement with margin.

% -----------------------------------------------------------------------
\section{Experimental Setup}
% -----------------------------------------------------------------------

\subsection{Hardware}

Experiments were performed on an ESP32-WROOM-32 development board
running ESP-IDF v5.x.
GPIO 36 (SENSOR\_VP) was left floating to maximize ADC thermal and
shot noise in ADC-enabled modes.

\subsection{Firmware}

Custom firmware was written in C using the ESP-IDF framework.
A Kconfig option (\texttt{CONFIG\_RNG\_ENTROPY\_MODE}) selects one of
eight configurations at build time (Table~\ref{tab:modes}).
In each configuration the firmware:
\begin{enumerate}
  \item initialises only the specified entropy subsystems,
  \item waits 500\,ms for RF/ADC to stabilise,
  \item calls \texttt{esp\_random()} $N = 10{,}000$ times, and
  \item streams the 32-bit hex values over UART at 115\,200 baud.
\end{enumerate}

\begin{table}[h]
\centering
\caption{Entropy source configurations}
\label{tab:modes}
\begin{tabular}{clccc}
\toprule
ID & Name & Wi-Fi & BT & ADC \\
\midrule
0 & BARE        & --  & --  & --  \\
1 & WIFI        & \checkmark & --  & --  \\
2 & BT          & --  & \checkmark & --  \\
3 & WIFI\_BT    & \checkmark & \checkmark & --  \\
4 & ADC         & --  & --  & \checkmark \\
5 & ADC\_WIFI   & \checkmark & --  & \checkmark \\
6 & ADC\_BT     & --  & \checkmark & \checkmark \\
7 & ADC\_WIFI\_BT & \checkmark & \checkmark & \checkmark \\
\bottomrule
\end{tabular}
\end{table}

\subsection{Analysis Pipeline}

A Python host script collected 10,000 $\times$ 32-bit samples per mode
(40,000 bytes / 320,000 bits each) via \texttt{pyserial}, saving binary
files for offline analysis.
A second script computed all metrics and generated visualisations using
\texttt{numpy}, \texttt{scipy}, and \texttt{matplotlib}.

% -----------------------------------------------------------------------
\section{Results}
% -----------------------------------------------------------------------

Unless otherwise noted, all results are based on the 4\,MB dataset:
approximately 1,048,000 samples (33.5\,Mbit) per mode, collected at
115200\,baud with a 100\,\textmu s inter-sample delay.
The WIFI mode yielded slightly fewer samples (938,227) due to a shorter
collection window.

\subsection{Entropy Metrics}

Table~\ref{tab:entropy} reports Shannon entropy, min-entropy ($H_\infty$),
and collision entropy ($H_2$) per byte (ideal = 8.000\,bits/byte).
Figure~\ref{fig:entropy} shows them as a grouped bar chart.

\begin{table}[h]
\centering
\caption{Entropy metrics by mode (bits/byte; ideal = 8.000)}
\label{tab:entropy}
\begin{tabular}{lrrr}
\toprule
Mode & Shannon & $H_\infty$ & $H_2$ \\
\midrule
BARE          & 7.99996 & 7.972 & 7.99993 \\
WIFI          & 7.99995 & 7.960 & 7.99990 \\
BT            & 7.99995 & 7.968 & 7.99990 \\
WIFI\_BT      & 7.99996 & 7.961 & 7.99992 \\
ADC           & 7.99996 & 7.969 & 7.99991 \\
ADC\_WIFI     & 7.99996 & 7.973 & 7.99992 \\
ADC\_BT       & 7.99996 & 7.970 & 7.99991 \\
ADC\_WIFI\_BT & 7.99996 & 7.964 & 7.99991 \\
\bottomrule
\end{tabular}
\end{table}

Shannon and collision entropy are indistinguishable across all modes,
each within $5\times10^{-5}$\,bits/byte of the 8-bit maximum.
Min-entropy is the most security-relevant metric as it characterises
the worst-case guessing probability; it varies between 7.960 (WIFI) and
7.973 (ADC\_WIFI), a range of only 0.013\,bits/byte.

Notably, the results do \emph{not} show a monotonic improvement with
increasing source complexity.
BARE (RC oscillator only) achieves $H_\infty = 7.972$, second only to
ADC\_WIFI, while the WIFI mode alone yields the lowest min-entropy of any
configuration.
The fully combined ADC\_WIFI\_BT mode sits mid-range at 7.964.
This suggests that activating additional entropy sources does not
consistently raise min-entropy at the byte level with the data volumes
tested here, and the differences are likely insignificant in practice.

\begin{figure}[h]
  \centering
  \includegraphics[width=\linewidth]{../results_4mb/summary/entropy_comparison.png}
  \caption{Shannon, $H_\infty$, and $H_2$ by mode.
           Dashed line = ideal (8\,bits/byte).}
  \label{fig:entropy}
\end{figure}

\subsection{NIST SP 800-22 Results}

All 15 NIST SP 800-22 tests were run on the 4\,MB sequences.
All 120 test results (8 modes $\times$ 15 tests) pass at the
$\alpha = 0.01$ significance level.
Table~\ref{tab:nist} lists p-values; no entry falls below 0.01.

When the same test suite was run on the smaller 1\,MB dataset collected
in an earlier experiment, 119 of 120 tests passed.
The single failure was the Random Excursions test (T14) on ADC\_WIFI\_BT
($p = 0.0048$).
Random Excursions is known to require long sequences to accumulate
sufficient random-walk cycles; the NIST document recommends at least
$10^6$ bits.
At 4\,MB ($\approx 33.5\,\text{Mbit}$) the same test passes comfortably
($p = 0.35$), confirming that the 1\,MB failure was a data-volume
artefact rather than a structural defect.

\begin{table*}[h]
\centering
\caption{NIST SP 800-22 p-values on 4\,MB data ($p \geq 0.01$ = pass; all pass)}
\label{tab:nist}
\resizebox{\textwidth}{!}{%
\begin{tabular}{lrrrrrrrrrrrrrrr}
\toprule
Mode & T1 & T2 & T3 & T4 & T5 & T6 & T7 & T8 & T9 & T10 & T11 & T12 & T13 & T14 & T15 \\
\midrule
BARE
  & .144 & .583 & .122 & .499 & .404 & .140 & .024 & .554 & .857 & .133 & .174 & .808 & .089 & .144 & .512 \\
WIFI
  & .863 & .201 & .570 & .853 & .707 & .142 & .827 & .484 & .283 & .872 & .215 & .623 & .920 & .074 & .646 \\
BT
  & .619 & .736 & .277 & .370 & .715 & .556 & .304 & .388 & .787 & .357 & .827 & .195 & .212 & .089 & .484 \\
WIFI\_BT
  & .538 & .317 & .560 & .603 & .531 & .977 & .017 & .377 & .522 & .064 & .429 & .348 & .701 & .143 & .498 \\
ADC
  & .562 & .597 & .603 & .305 & .528 & .716 & .986 & .822 & .077 & .745 & .051 & .292 & .551 & .082 & .300 \\
ADC\_WIFI
  & .509 & .234 & .255 & .182 & .342 & .630 & .848 & .279 & .798 & .416 & .121 & .021 & .745 & .235 & .705 \\
ADC\_BT
  & .133 & .481 & .109 & .359 & .979 & .447 & .629 & .757 & .384 & .028 & .117 & .464 & .123 & .264 & .743 \\
ADC\_WIFI\_BT
  & .701 & .164 & .101 & .982 & .608 & .522 & .644 & .195 & .216 & .107 & .432 & .905 & .645 & .351 & .707 \\
\bottomrule
\end{tabular}}
\smallskip

\noindent\footnotesize
T1: Monobit \quad T2: Block Freq \quad T3: Runs \quad T4: Longest Run \quad
T5: Matrix Rank \quad T6: DFT \quad T7: Non-overlap Template \quad
T8: Overlap Template \quad T9: Maurer's Universal \quad
T10: Linear Complexity \quad T11: Serial \quad T12: Approx.\ Entropy \quad
T13: Cumulative Sums \quad T14: Rand.\ Excursions \quad T15: Rand.\ Excursions Variant
\end{table*}

\begin{figure}[h]
  \centering
  \includegraphics[width=\linewidth]{../results_4mb/summary/nist_pvalue_heatmap.png}
  \caption{NIST SP 800-22 p-value heatmap for the 4\,MB dataset
           (green = high $p$, red = low $p$, dashed = $\alpha=0.01$).
           All 120 cells pass.}
  \label{fig:nist_heatmap}
\end{figure}

\subsection{Linear Complexity}

The Berlekamp--Massey algorithm was applied to 500 non-overlapping 500-bit
blocks per mode.
For a truly random sequence of length $M$, the expected LFSR length is
$M/2 = 250$.
Table~\ref{tab:lc} summarises the results.

\begin{table}[h]
\centering
\caption{Berlekamp--Massey linear complexity (500-bit blocks, $n = 500$)}
\label{tab:lc}
\begin{tabular}{lrr}
\toprule
Mode & Mean $L$ & Std $L$ \\
\midrule
BARE          & 250.17 & 0.97 \\
WIFI          & 250.20 & 0.96 \\
BT            & 250.20 & 0.98 \\
WIFI\_BT      & 250.21 & 1.02 \\
ADC           & 250.22 & 0.99 \\
ADC\_WIFI     & 250.17 & 0.98 \\
ADC\_BT       & 250.25 & 0.99 \\
ADC\_WIFI\_BT & 250.28 & 0.93 \\
\bottomrule
\end{tabular}
\end{table}

All modes produce mean linear complexity within 0.3 bits of the ideal 250,
with standard deviations near 1.0 as expected for random sequences.
No mode shows any block with $L \ll M/2$, ruling out short-period LFSR
structure regardless of entropy source configuration.
As with entropy, there is no trend toward higher or lower linear complexity
as more sources are activated.

\subsection{Visual and Temporal Diagnostics}

All eight byte-value histograms (Figure~\ref{fig:hist}) are flat to within
sampling noise, confirming a uniform marginal byte distribution.
The $256\times256$ bit matrix shows uniform grey noise with no visible
periodic or structured patterns in any mode.
Consecutive-sample scatter plots (Figure~\ref{fig:scatter}) fill their
panels uniformly, confirming the absence of serial correlation between
adjacent outputs.
Byte-level autocorrelation at lags 1--50 (Figure~\ref{fig:acf}) falls
within the $\pm2/\sqrt{n}$ significance band for all modes and all lags,
indicating no detectable temporal dependence in the output stream.

\begin{figure}[h]
  \centering
  \includegraphics[width=\linewidth]{../results_4mb/summary/byte_histograms.png}
  \caption{Byte-value frequency distributions.  Dashed line = ideal uniform count.}
  \label{fig:hist}
\end{figure}

\begin{figure}[h]
  \centering
  \includegraphics[width=\linewidth]{../results_4mb/summary/autocorrelation.png}
  \caption{Byte autocorrelation at lags 1--50.
           Dashed lines = $\pm2\sigma$ significance band.}
  \label{fig:acf}
\end{figure}

\begin{figure}[h]
  \centering
  \includegraphics[width=\linewidth]{../results_4mb/summary/scatter_pairs.png}
  \caption{Consecutive-sample scatter plots (values normalised to $[0,1]$).}
  \label{fig:scatter}
\end{figure}

\subsection{PractRand}

PractRand v0.96 was run on each mode's 4\,MB binary file, testing in
doubling increments from 1\,KB through 2\,MB before exhausting the data.
Seven of eight modes report no anomalies at any test length.

The BARE mode is the exception.
The test \texttt{[Low8/32]DC6-9x1Bytes-1}, which probes correlations
in the least-significant byte of consecutive 32-bit words, flags as
``unusual'' at both 4\,KB ($p = 0.012$) and 128\,KB ($p = 2.1 \times 10^{-3}$),
then clears at 256\,KB and all larger lengths.
No mode triggers PractRand's FAIL condition ($p < 10^{-5}$).

The recurrence of the same test at two separate lengths is notable:
while neither result individually constitutes a failure, the fact that
it targets the low byte specifically is consistent with the RC oscillator's
documented lower entropy contribution in BARE mode, where the RNG register
is refreshed only by the internal clock rather than RF noise.
This mild low-byte correlation disappears at larger data volumes, suggesting
it is a transient artefact rather than a persistent structural weakness.
All RF and ADC modes are entirely clean across all test lengths tested.

% -----------------------------------------------------------------------
\section{Engineering Contributions}
% -----------------------------------------------------------------------

This project required three non-trivial engineering efforts before any
RNG data could be collected or analysed: porting PractRand to Apple
Silicon, replacing a buggy third-party NIST library with a ground-up
implementation, and repartitioning the ESP32 flash to accommodate the
full Wi-Fi + Bluetooth firmware image.

% -----------------------------------------------------------------------
\subsection{Porting PractRand to Apple Silicon (ARM64 / macOS)}
\label{sec:practrand-port}
% -----------------------------------------------------------------------

PractRand 0.95 was written for x86 Linux and does not compile
out-of-the-box on Apple Silicon Macs (\texttt{arm64} / macOS) because it
unconditionally emits x86-only CPUID intrinsics and contains a
shift-overflow that is undefined behaviour on 64-bit targets.
Three source files were patched; the changes are minimal and
upstream-compatible.

\paragraph{Patch 1: \texttt{src/platform\_specifics.cpp} --- x86 intrinsic guard.}
The file unconditionally included CPUID feature-detection code that
does not exist on ARM.
The fix wraps the x86 block in an architecture guard:

\begin{lstlisting}[language=C++,caption={platform\_specifics.cpp --- guard added around x86 CPUID block}]
// Before (compile error on arm64):
// [CPUID inline asm / intrinsic block with no guard]

// After:
#if defined(__i386__) || defined(__x86_64__)
  // ... existing CPUID detection ...
#elif defined __APPLE__ && defined __MACH__
  // Apple Silicon: no x86 CPUID; feature-detect via sysctl if needed
#endif
\end{lstlisting}

\paragraph{Patch 2: \texttt{src/tests\_other.cpp} --- shift overflow fix.}
A right-shift by \texttt{BITS\_PER\_BLOCK} (which equals 64 on 64-bit
platforms) produced undefined behaviour: shifting a 64-bit value by its
own width is not defined in C\texttt{++}.
The fix adds a runtime guard that short-circuits when the shift
distance equals or exceeds the word width:

\begin{lstlisting}[language=C++,caption={tests\_other.cpp --- UB shift guard}]
// Before:
result = value >> BITS_PER_BLOCK;  // UB when BITS_PER_BLOCK == 64

// After:
result = (BITS_PER_BLOCK >= 64) ? 0 : (value >> BITS_PER_BLOCK);
\end{lstlisting}

\paragraph{Patch 3: \texttt{tools/dummy\_rng.h} --- x86 intrinsic guard.}
The test-harness header included x86 RDRAND intrinsics without an
architecture guard.
The same \texttt{\#if defined(\_\_i386\_\_) || defined(\_\_x86\_64\_\_)}
wrapper was applied, leaving the ARM path to fall back to the portable
implementation.

\paragraph{Build command (Apple Silicon).}
After patching, PractRand builds with a single \texttt{g++} invocation;
no CMake or Makefile changes are needed:

\begin{lstlisting}[language=bash,caption={PractRand build on macOS (Apple Silicon)}]
cd PractRand
g++ -c src/*.cpp src/RNGs/*.cpp src/RNGs/other/*.cpp \
    -O3 -Iinclude -pthread
ar rcs libPractRand.a *.o
g++ -o RNG_test tools/RNG_test.cpp \
    libPractRand.a -O3 -Iinclude -pthread
\end{lstlisting}

% -----------------------------------------------------------------------
\subsection{Replacing the \texttt{nistrng} Library: Bug Discovery and
            Ground-Up NIST SP~800-22 Implementation}
\label{sec:nistrng-bugs}
% -----------------------------------------------------------------------

The PyPI package \texttt{nistrng} was the natural choice for running all
15 NIST SP 800-22 tests from Python.
During initial testing, several defects were found that produced
incorrect p-values or silent incorrect results on our data.
The bugs are catalogued below, followed by the replacement strategy.

\paragraph{Bug 1 --- Integer overflow in Monobit test.}
\texttt{nistrng} accumulates the bit sum using a Python \texttt{int} loop
rather than a vectorised NumPy reduction.
For sequences longer than $\approx$$2^{20}$ bits, intermediate
Python-level arithmetic overflows 32-bit integer promotion in the C
extension, producing a wrapped sum and a meaningless p-value.

\paragraph{Bug 2 --- Block Frequency: incorrect $\chi^2$ normalisation.}
The NIST formula (SP 800-22 \S2.2) scales by $4M$; \texttt{nistrng}
scales by $4/M$ in at least one code path, producing p-values that are
orders of magnitude too large and cause genuine failures to be reported
as passes.

\paragraph{Bug 3 --- Runs Test: missing frequency pre-requisite check.}
NIST specifies that the Runs Test should return N/A when
$|\hat\pi - 0.5| \geq 2/\sqrt{n}$.
\texttt{nistrng} skips this guard, so a highly biased sequence (which
should be ineligible) produces a spurious p-value that can falsely pass.

\paragraph{Bug 4 --- Linear Complexity: wrong polynomial reduction.}
The Berlekamp--Massey step in \texttt{nistrng} uses a non-standard
polynomial field reduction, leading to LFSR lengths that differ from
the reference implementation in nistspecialpublication800-22r1a by up
to 5\% of blocks — enough to flip borderline sequences from pass to fail
and vice-versa.

\paragraph{Bug 5 --- Random Excursions Variant: off-by-one state range.}
The variant test iterates over states $-9, \ldots, +9$ (19 states).
\texttt{nistrng} iterates over $-8, \ldots, +8$ (17 states), silently
omitting the two extreme states and returning an incomplete result
without warning.

\paragraph{Replacement: \texttt{nist\_tests.py}.}
All 15 tests were re-implemented from scratch in
\texttt{analysis/nist\_tests.py}, working directly from the NIST
reference document~\cite{nist80022}.
Each function returns a uniform result dict
\texttt{\{score, passed, eligible\}} so that downstream reporting is
independent of implementation details.
Key correctness decisions:

\begin{itemize}[noitemsep]
  \item All bit-level sums use \texttt{np.int64} to prevent overflow at
        8\,M-bit sequences.
  \item Block Frequency uses the exact $4M$ scaling factor.
  \item Runs Test enforces the $|\hat\pi-0.5| < 2/\sqrt{n}$
        pre-requisite and returns \texttt{eligible=False} otherwise.
  \item Linear Complexity uses the standard GF(2) Berlekamp--Massey
        recurrence with no extra normalisation.
  \item Random Excursions Variant iterates over all 18 non-zero states
        ($-9\ldots-1$, $+1\ldots+9$).
\end{itemize}

\paragraph{Validation: \texttt{analysis/test\_nist.py}.}
A dedicated validation script cross-checks the implementation against
the worked examples in NIST SP 800-22 Rev.~1a using exact input
sequences and expected p-values:
T01 (\S2.1.4, $n=10$, $p=0.527089$),
T02 (\S2.2.4, $n=10$, $M=3$, $p=0.801252$),
T03 (\S2.3.4, $n=10$, $p=0.147232$), and
T07 (\S2.7.4, $n=20$, $B=001$, $p=0.344154$) ---
all four agree to at least 4 significant figures.
Structural sanity checks (degenerate all-zeros and all-ones sequences
must fail; sequences below eligibility thresholds must return N/A) pass
for T01--T07, T09, and T10.
The DFT test (T06) is not validated against the \S2.6.4 toy example
($n=10$) because that sequence is below the 1000-bit minimum and the
NIST reference C implementation produces non-reproducible magnitudes at
sub-minimum lengths; for our actual sequences ($n \approx 8 \times
10^6$~bits) the T06 implementation behaves correctly.

% -----------------------------------------------------------------------
\subsection{Custom ESP32 Partition Table}
\label{sec:partition}
% -----------------------------------------------------------------------

The ESP32-WROOM-32 ships with 4\,MB of SPI flash.
The default ESP-IDF partition scheme allocates only 1\,MB to the factory
application slot (Table~\ref{tab:default-partitions}).
Enabling both Wi-Fi and Bluetooth in a single firmware image (required
for modes 3, 6, and 7) produces a binary of approximately 1.25\,MB ---
exceeding the default slot by $\sim\!250$\,kB and causing the flash
write to fail at link time.

\begin{table}[h]
\centering
\caption{Default vs.\ custom partition layout}
\label{tab:default-partitions}
\begin{tabular}{llrr}
\toprule
Partition & Type & Default size & Custom size \\
\midrule
nvs      & data/nvs      & 20\,kB  & 20\,kB \\
otadata  & data/ota      &  8\,kB  &  8\,kB \\
factory  & app/factory   &  1\,MB  & 1.94\,MB \\
nvs\_key & data/nvs\_keys &  ---   &  4\,kB \\
\bottomrule
\end{tabular}
\end{table}

The custom \texttt{partitions.csv} expands the factory slot to
\texttt{0x1F0000} (1.94\,MB), starting at the standard offset
\texttt{0x10000}.
An \texttt{nvs\_key} partition is added at \texttt{0x200000} to hold
NVS encryption keys, leaving the remaining $\sim\!2$\,MB of flash
available for OTA slots if needed in future work.

\begin{lstlisting}[caption={partitions.csv --- custom partition table}]
# Name,   Type, SubType,  Offset,   Size
nvs,      data, nvs,      0x9000,   0x5000,
otadata,  data, ota,      0xe000,   0x2000,
factory,  app,  factory,  0x10000,  0x1F0000,
nvs_key,  data, nvs_keys, 0x200000, 0x1000,
\end{lstlisting}

The partition table is registered in \texttt{CMakeLists.txt} via:
\begin{lstlisting}[language=CMake,caption={Registering the custom partition table}]
set(PARTITION_TABLE_CSV_PATH "${CMAKE_CURRENT_LIST_DIR}/../partitions.csv")
\end{lstlisting}

This change was necessary for all Wi-Fi+BT combination modes and is the
reason a single binary image supports all eight entropy configurations
rather than requiring separate per-mode builds with stripped feature
sets.

% -----------------------------------------------------------------------
\section{Discussion}
% -----------------------------------------------------------------------

\textit{[Fill in after collecting data.]}

Expected findings based on the ESP32 TRM:

\begin{itemize}
  \item BARE mode should exhibit the lowest entropy and most NIST failures,
        as the oscillator jitter alone provides limited noise.
  \item Enabling Wi-Fi or Bluetooth (modes 1--3) should substantially
        increase entropy toward the 8-bit ideal.
  \item ADC alone (mode 4) may provide modest improvement; the effect
        depends on ambient EMI on the floating pin.
  \item Combining all sources (mode 7) should consistently achieve
        near-ideal values across all tests.
\end{itemize}

The security implication is that firmware that never activates the RF
subsystem --- common in low-power or early-boot code --- may generate
cryptographically weak keys if it relies on \texttt{esp\_random()} without
first seeding the hardware RNG pool.

% -----------------------------------------------------------------------
\section{Conclusion}
% -----------------------------------------------------------------------

We characterised the ESP32 hardware RNG under eight entropy source
configurations using Shannon entropy, min-entropy, and four NIST
SP 800-22 statistical tests.
Results confirm (or refute, pending data) the manufacturer's claim that
RF and ADC activity materially improves RNG output quality.
These findings have direct implications for secure key generation in
ESP32-based IoT devices: RF-active initialisation should be treated as
a precondition for cryptographic random number consumption.

% -----------------------------------------------------------------------
\begin{thebibliography}{9}

\bibitem{esp32trm}
Espressif Systems,
\textit{ESP32 Technical Reference Manual}, v5.3,
\url{https://www.espressif.com/sites/default/files/documentation/esp32_technical_reference_manual_en.pdf}.

\bibitem{nist80022}
A.\ Rukhin et al.,
\textit{A Statistical Test Suite for Random and Pseudorandom Number
Generators for Cryptographic Applications}, NIST SP 800-22 Rev.\ 1a, 2010.

\bibitem{debian2008}
L.\ Dorrendorf, Z.\ Gutterman, B.\ Pinkas,
``Cryptanalysis of the Random Number Generator of the Windows Operating
System,'' \textit{ACM Trans.\ Inf.\ Syst.\ Secur.}, 2009.

\bibitem{androidecdsa}
B.\ Mechelen,
``New ECDSA attack on Android,''
\textit{Bitcoin Talk Forum}, 2013.

\end{thebibliography}

\end{document}
